% English translation of the verified Chinese baseline. Structural identifiers must remain synchronized.
\section{Discussion and Scope of Applicability}\label{sec:discussion}

\subsection{Theoretical foundations and methodological boundaries}\label{sec:discussion-7-1}

The continuum mechanics and finite-dimensional reduction are founded on geometrically exact rod theory, classical geometric mechanics, and geometric numerical integration. The HNN, dissipative/port-network, and complex-valued-network literature defines the learnable objects and function classes. Finite-dimensional continuous KAM theory, KAM theory for symplectic maps, and the parameterization method support the conservative flow, symplectic discretization, and invariant-torus validation, respectively. These bodies of literature concern different mathematical objects; software documentation or learning-model literature is not used as a substitute for the structural assumptions of continuum mechanics or KAM theorems.\cite{Simo1985FiniteStrainBeam,Antman2005Elasticity,MarsdenRatiu1999MechanicsSymmetry,HairerLubichWanner2006GNI,Greydanus2019HNN,Zhong2020Dissipative,Desai2021PHNN,Chen2022KAMSLT,Trabelsi2018DeepComplex,Hirose2012CVNN,FengQin2010,Shang1999KAMAlgorithms,Shang2000KAMMappings,Haro2016Parameterization,Arnold1989ClassicalMechanics,Poschel2001KAMLecture}

Moser's invariant-curve theorem for annulus maps, Kuksin's infinite-dimensional KAM theory, the metriplectic framework, and theoretical results for specific complex-reaction networks delimit the applicability of two-dimensional sections, PDE limits, thermodynamic dissipation, and architecture-specific advantages of complex structures. They do not directly constitute the principal theorems for the finite-dimensional catheter system considered here.\cite{Coquinot2020Metriplectic,Moser1962InvariantCurves,Kuksin1993InfiniteDimensionalKAM,Poschel2001KAMLecture,Zhang2022ComplexReaction}

The canonical data-contract diagnosis shows that structural consistency must be checked componentwise in the configuration, tangent, and cotangent spaces and cannot be replaced by energy drift or a single residual. Cross-trajectory comparison, the distinction between projected momentum and the mass diagnostic, the unavailability of an internal solver stage, and omission of the fixed-basis transform from initialization represent four distinct error sources: comparison objects, cotangent projection, output semantics, and initialization. The two-level time-step study did not meet the predefined convergence criterion, whereas the three-level study resolved the convergence trend. This sequence provides a reusable diagnostic strategy for other finite-dimensional projection problems. Force-v2 closure and Bridge perturbation-energy drift address different questions and cannot substitute for one another.

\subsection{Present scope and conditions requiring further validation}\label{sec:discussion-7-2}

The continuous model, joint Galerkin reduction, canonical Hamiltonian formulation, structured CHNN, error budget, and invariant-torus validation interfaces have been constructed, and the COMSOL field projection and canonical-label mapping have been instantiated on a fixed \(N_{\mathrm{diag}}=2\) basis. P0, P1, P1.5, and P1.6 provide low-dimensional and three-trajectory small-sample results; P1.6 specifically shows that exact satisfaction of hard structural constraints can coexist with degraded prediction. The three-level and \(1\,\mathrm{s}\) studies support at most this single diagnostic sentinel as a candidate input to a subsequent A1 qualification experiment; they do not certify data-generation qualification across modes, amplitudes, or operating conditions.

The principal modal order \(n_{\mathrm{main}}\), full 96-trajectory data set, P2 model comparison, open-port identification, high-fidelity discrepancy, normal-form coefficients, nondegeneracy constants, statistical-learning constants, high-order error certificates, and KAM thresholds still require numerical instantiation. Seeds in the P1 series are algorithmic repetitions, and three trajectories do not support population-level physical extrapolation. The present results therefore do not establish general CHNN superiority over Real HNN on the full P2 task, invariant-torus certification for a catheter operating condition, a long-time stability guarantee for the complete open system, an engineering safe domain, or clinical efficacy. BDF remains only a state-sampling and vector-field-labeling method and has not been shown to be symplectic.

The code-material boundary is equally explicit: the current repository contains only \texttt{verification\_skeleton.py} and \texttt{verification\_results.json}. They are an interface example, regression-test skeleton, numerical-workflow demonstration, and build-gate test object with its result, not formal CHNN, a formal COMSOL--Galerkin bridge, formal force-v2 code, a P0/P1/P2 production experiment, a controller, an adaptive law, or a numerical KAM certificate. Removing this boundary requires formal source code, configuration, an environment lock, raw arrays, a manifest, clean-run logs, automatically generated tables and figures, and provenance-consistency records; while those materials are missing, A-17 remains PARTIAL.

\subsection{Auxiliary models and future extensions}\label{sec:discussion-7-3}

A strict weak-form COMSOL implementation, Timoshenko beams, and three-dimensional solid models will support high-fidelity validation. Infinite-dimensional KAM theory, thermodynamic metriplectic structures, learning of unknown canonical coordinates, complete Maxwell-field models, and moving-wall fluid--structure interaction are retained as future extensions. Open-port studies will further examine power balance, dissipation sign, and port identifiability, while numerical KAM studies will compute nondegeneracy margins, high-order errors, and invariant-torus residuals. These models and validation tasks retain their full scientific scope but are not conflated with the present conclusions for the finite-dimensional conservative Kirchhoff branch.\cite{COMSOLBeamInterface2025,COMSOLGeometricNonlinearity2025,COMSOLSectionStiffness2025,COMSOLWeakFormPDE2025,COMSOLWeakContribution2025,Jackson1998Electrodynamics,Abbott2020MagneticRobotics,Coquinot2020Metriplectic,Kuksin1993InfiniteDimensionalKAM}
